\section{Retos de los espacios de acción continuos}

La forma original de DQN solo funciona con espacios de acción \textbf{discretos} (la red produce el valor Q de cada acción y luego se toma el argmax). Para espacios de acción continuos, $\max_a Q(s,a)$ es en sí mismo un problema de optimización.

Las soluciones incluyen:
\begin{enumerate}
    \item \textbf{Discretización}: discretizar el espacio continuo, aunque la maldición de la dimensionalidad es severa.
    \item \textbf{Muestreo}: muestrear aleatoriamente varias acciones y tomar la de mayor valor Q.
    \item \textbf{Aprender un actor}: volver al marco de actor-crítico (como SAC).
    \item \textbf{DDPG}: entrenar una red de política determinista $\mu_\theta(s) = \arg\max_a Q(s,a)$.
\end{enumerate}

\begin{warningbox}{Dificultades de Q-Learning en espacios continuos}
Q-Learning puro en un espacio de acción continuo necesita resolver el problema de optimización interno $\max_a Q(s,a)$, lo cual es computacionalmente costoso en espacios de alta dimensión. Esta es también la razón por la que, en tareas de control continuo, los métodos de actor-crítico (como SAC) suelen preferirse frente a Q-learning puro.
\end{warningbox}

\subsection{Resumen de la sección}
 
Q-Learning es natural y eficiente en espacios de acción discretos, pero en espacios de acción continuos requiere técnicas adicionales o volver a los métodos de actor-crítico.

